\scanpage{302}
\mspage{452}

Let us note that (33), can be written $\rho=-(1+N_1)\sigma$ i.e.
\deleted{\illegible{1}}
\[
\rho+\sigma=-N_1\sigma \tag{35}
\]
multiplying by $N_1$ on the left we get
\[
N_1\rho+N_1\sigma=0
\]
(since $N_1^2=0$), moreover
\[
N_1\rho=\rho N_0,\qquad N_1\sigma=\sigma N_0
\]
(for $N_1=\rho N_0\rho^{-1}=\sigma N_0\sigma^{-1}$), so (\unsure{35}) can be written
\[
\rho+\sigma=-N_1\sigma=+N_1\rho=-\sigma N_0=+\rho N_0
\overset{\text{definition}}{=}N . \tag{\unsure{35|36}}
\]
moreover when $\rho$ is in the general form
$\rho=\begin{pmatrix}0&u\\v&w\end{pmatrix}$ we find
\[
N\deleted{\illegible{1}}=\begin{pmatrix}0&0\\0&u\end{pmatrix},
\qquad N^2=\begin{pmatrix}0&0\\0&u^2\end{pmatrix}
\]
so the normalization (34) $u=1$ also means
\[
\rho+\sigma=\begin{pmatrix}0&0\\0&1\end{pmatrix}\deleted{\unsure{by}}, \tag{37}
\]
\inserted{or again (since $u$ is invertible)
\[
N^2=N\qquad (\text{where }N=\rho+\sigma) \tag{37 bis}
\]}
or again that $\sigma$ is given in terms of $\rho$ by
\[
\sigma=\begin{pmatrix}0&-1\\-v&1-w\end{pmatrix}. \tag{38}
\]
moreover, $N_0$ \inserted{$=\rho N_1\rho^{-1}$} will be of the form
$\begin{pmatrix}0&\alpha\\0&0\end{pmatrix}$,
an immediate calculation shows that $\alpha$ is given by the formula
$\unsure{u}=\alpha v$, i.e. (\unsure{if} we normalize by $u=1$)

\margin{\uline{N.B.} We see \unsure{at once} that the conditions considered so far
\unsure{imply} that $1,\rho,\sigma$ \unsure{are linearly independent}
\unsure{in each} fiber, i.e. $T_\rho\cap T_\sigma=\unsure{\mathbb G_m}$ !
i.e. $\widetilde T_\rho\cap\widetilde T_\sigma=\{1\}$.}

\scanpage{303}

\[
N_0=\begin{pmatrix}0&1/v\\0&0\end{pmatrix},\qquad
\varepsilon_0=\begin{pmatrix}1&1/v\\0&1\end{pmatrix}. \tag{39}
\]

\margin{\uline{N.B.} We will have
\[
\boxed{\sigma^{-1}\rho=-\varepsilon_0}\deleted{\unsure{N_0}} \tag{39 bis}
\]
For $\sigma^{-1}\rho=\sigma^{-1}(\rho\sigma^{-1})\sigma
=\sigma^{-1}(-\varepsilon_1)\sigma
=-\sigma^{-1}(\varepsilon_1)=-\varepsilon_0$.}

\deleted{Let us return}
These canonical normalizations having been made, let us return to the proposition on page 436 ; since $\varepsilon_1$ has been chosen
(in as canonical a way as can be) there is no longer any drawback to
\deleted{\illegible{1}}\inserted{leaving} equation (32) (put into the apparently more intrinsic form
(11) \unsure{or} (12)) in the \deleted{\illegible{1}} initial form,
where a priori we were thinking of introducing the parameter $\gamma$ with
\[
[\alpha,\rho]=[\beta,\rho]\varepsilon_1(\gamma)\unsure{\,;}\quad -
\]
but it appears that it is $\gamma+1=\mu$ that has a particularly remarkable significance
in the theory --- so we \unsure{keep|will keep} its form (32), with the factor $\varepsilon_1(\mu-1)$.

The expression (97) for $P(\unsure{\zeta})$, thanks to (37), can be written
\[
P(\zeta)=D\qquad\text{for }\zeta=\begin{pmatrix}A&B\\C&D\end{pmatrix}, \tag{40}
\]

\scanpage{304}
\mspage{453}

Thus the proposition tells us that the scheme \inserted{$\mathcal X_{\rho,\sigma}$}
of simultaneous solutions in $(\zeta,\eta,\mu)$ of
\[
\det\zeta=\det\eta=1,\qquad \zeta=\eta\varepsilon_1(\mu-1),\qquad
\operatorname{Tr}(\zeta\rho)=\operatorname{Tr}\rho,\qquad
\operatorname{Tr}\eta\deleted{\unsure{\zeta|\rho}}\inserted{\sigma}
=\operatorname{Tr}\unsure{\sigma|\alpha} \tag{41}
\]
is isomorphic, via $(\zeta,\eta,\mu)\longrightarrow\zeta$, to that of the
$\zeta=\begin{pmatrix}A&B\\C&D\end{pmatrix}$ satisfying the conditions
\[
\underbrace{\det\zeta}_{AD-BC}=1,\qquad
\underbrace{\operatorname{Tr}\zeta\rho}_{(C+wD)+vB}
=\underbrace{\operatorname{Tr}\rho}_{w},\qquad
\text{and }D\text{ an invertible section of }\mathcal O_S \tag{42}
\]
and we recover $\mu$ in terms of $\zeta$ by
\[
\mu=1/D \tag{43}
\]
(here $c=\operatorname{Tr}(\rho+\sigma)=\operatorname{Tr}\begin{pmatrix}0&0\\0&1\end{pmatrix}=1$),
whence $\eta$ by $\eta=\zeta\varepsilon_1(1-\mu)$.

The subgroup scheme $B'_0$ of page 442$'$ is none other than the scheme of matrices
\[
B'_0=\left\{\begin{pmatrix}\mu&\lambda\\0&1\end{pmatrix}
\,\middle|\,\mu\in\Gamma\mathbb G_m,\ \lambda\in\Gamma\mathbb G_a\right\} \tag{44}
\]
We are going to make explicit the isomorphism (79)
\[
\underline u\longmapsto\bigl([\underline u,\rho],[\underline u,\sigma],\det\underline u\bigr) \tag{45}
\]
in terms of the coefficients $\mu,\lambda$ of $B'_0$, with the help of the parameters
$v,w$ in $\rho=\begin{pmatrix}0&1\\v&w\end{pmatrix}$.

\scanpage{305}

We find
\deleted{\unsure{$\underline u(\rho)=$}}
\[
\begin{aligned}
{}[\underline u,\rho]
&=\begin{pmatrix}\mu&\lambda\\0&1\end{pmatrix}
\begin{pmatrix}0&1\\v&w\end{pmatrix}
\frac1\mu\begin{pmatrix}1&-\lambda\\0&\mu\end{pmatrix}
\left(-\frac1v\right)\begin{pmatrix}w&-1\\-v&0\end{pmatrix}\\
&=-\frac1{\mu v}\begin{pmatrix}\lambda v&\mu+\lambda w\\v&w\end{pmatrix}
\begin{pmatrix}w+\lambda v&-1\\-\mu v&0\end{pmatrix}\\
&=-\frac1{\mu v}
\begin{pmatrix}
\deleted{\illegible{2}\unsure{-\mu)vw+\lambda}}\inserted{v^2\lambda^2-v\mu^2-vw\lambda\mu+vw\lambda}&-\lambda v\\
v^2\lambda-vw\mu+vw&-v
\end{pmatrix}\\
&=\frac1\mu\begin{pmatrix}
-v\lambda^2+w\lambda\mu+\mu^2-w\lambda&\unsure{\lambda}\\
-v\lambda+w\mu-w&1
\end{pmatrix}
\end{aligned}
\]

\margin{\unsure{now|and} for $\sigma=\begin{pmatrix}0&-1\\v'&w'\end{pmatrix}$
with $v'=-v$, $w'=1-w$}

\[
\begin{aligned}
{}[\underline u,\sigma]
&=\begin{pmatrix}\mu&\lambda\\0&1\end{pmatrix}
\begin{pmatrix}0&-1\\v'&w'\end{pmatrix}
\frac1\mu\begin{pmatrix}1&-\lambda\\0&\mu\end{pmatrix}
\left(\frac1{v'}\right)\begin{pmatrix}w'&1\\-v'&0\end{pmatrix}\\
&=\frac1{\mu v'}\begin{pmatrix}\lambda v'&-\mu+\lambda w'\\v'&w'\end{pmatrix}
\begin{pmatrix}w'+\lambda v'&1\\-\mu w'&0\end{pmatrix}\\
&=\frac1{\mu v'}\begin{pmatrix}
v'^2\lambda^2+v'\mu^2-v'w'\lambda\mu+v'w'\lambda&\lambda v'\\
v'^2\lambda-v'w'\mu+v'w'&v'
\end{pmatrix}\\
&=\frac1\mu\begin{pmatrix}
v'\lambda^2\unsure{-}\deleted{\unsure{\mu}}w'\lambda\mu+\mu^2+w'\lambda&\lambda\\
v'\lambda-w'\mu+w'&1
\end{pmatrix}\\
&=\frac1\mu\begin{pmatrix}
-v\lambda^2+(w-1)\lambda\mu+\mu^2\unsure{-}(w-1)\lambda&\lambda\\
-v\lambda+(w-1)\mu-(w-1)\deleted{\unsure{\lambda}}&1
\end{pmatrix}
\end{aligned}
\]

In summary
\[
\left\{\begin{aligned}
\zeta=[\underline u,\rho]
&=\frac1\mu\begin{pmatrix}
\mu^2+(-v\lambda+w(\mu\unsure{-}1))\lambda&\lambda\\
-\deleted{\unsure{(}}v\lambda+w(\mu-1)&1
\end{pmatrix}\\
\eta=[\underline u,\sigma]
&=\frac1\mu\begin{pmatrix}
\mu^2+(-v\lambda+(w-1)(\mu-1))\lambda&\lambda\\
-v\lambda+(w-1)(\mu-1)&1
\end{pmatrix}\quad\text{and, of course}\\
\mu&=\det\underline u\ \deleted{\unsure{=\mu}}
\end{aligned}\right. \tag{46}
\]
so we have, if $\zeta=\begin{pmatrix}A&B\\C&D\end{pmatrix}$
\[
B=\lambda/\mu,\quad D=1/\mu\qquad\text{i.e.} \tag{47}
\]
\[
\mu=1/D,\quad\lambda=B/D,\qquad\text{\unsure{i.e.|of which} }\underline u=\begin{pmatrix}1/D&B/D\\0&1\end{pmatrix} \tag{48}
\]

\scanpage{306}
\mspage{453bis}

Likewise, if $\eta=\begin{pmatrix}R&S\\T&U\end{pmatrix}$, we will have \deleted{$S=B,U=D$}
\deleted{whence}
\[
\left\{\begin{aligned}
S=B&=\lambda/\mu,\quad U=D=1/\mu\qquad\text{whence}\\
\deleted{\unsure{and}}\ \mu&=1/U,\quad\lambda=S/U
\end{aligned}\right. \tag{47 bis}
\]
so
\[
\underline u=\begin{pmatrix}1/U\deleted{\unsure{(=\mu)}}&S/U\\0&1\end{pmatrix}. \tag{\unsure{48 bis}}
\]

Let us now consider the equation
\[
[\alpha,\rho]=[\beta,\sigma]\varepsilon_1(\mu-1),\qquad\text{with} \tag{49}
\]
\[
\alpha=\begin{pmatrix}a&b\\c&d\end{pmatrix},\quad
\beta=\begin{pmatrix}r&s\\t&u\end{pmatrix}, \tag{\unsure{50}}
\]
and set, for more \unsure{elegant} calculations
\[
\rho=\begin{pmatrix}0&\unsure{v_1}\\\unsure{v}&w\end{pmatrix},\quad
\sigma=\begin{pmatrix}0&v'_1\\v'&w'\end{pmatrix} \tag{51}
\]
(we have written $v_1,v'_1$ instead of $u,u'$, so as not to \unsure{interfere} with the notation $u$ for one of the coefficients of $\beta$), with
\[
v'=-v,\quad v'_1=-\unsure{v_1|v'},\quad w+w'=v_1,\qquad\text{let us set} \tag{52}
\]
\[
\delta_\alpha=\det\alpha=ad-bc,\quad
\delta_\beta=\det\beta=ru-st,
\tag{53}
\]

\margin{We do not set $v_1=1$, so as \unsure{not to have} to \unsure{redo} the same calculation for $[\beta,\sigma]$ as for $[\alpha,\rho]$ !}

we then find
\[
\begin{aligned}
\zeta=[\alpha,\rho]
&=\begin{pmatrix}a&b\\c&d\end{pmatrix}
\begin{pmatrix}0&v_1\\v&w\end{pmatrix}
\frac1{\delta_\alpha}\begin{pmatrix}d&-b\\-c&a\end{pmatrix}
\deleted{\unsure{(}}\frac1{-vv_1}\begin{pmatrix}w&-v_1\\-v&0\end{pmatrix}\\
&=\unsure{-}\frac1{\delta_\alpha vv_1}
\begin{pmatrix}bv&av_1+bw\\dv&cv_1+dw\end{pmatrix}
\begin{pmatrix}dw+bv&-dv_1\\-cw-av&cv_1\end{pmatrix}
=\begin{pmatrix}A&B\\C&D\end{pmatrix}
\end{aligned}
\]
with
\[
\left\{\begin{aligned}
(-\delta_\alpha vv_1)A&=\ldots\quad\text{(a bore).},
&(-\delta_\alpha vv_1)B&=-vv_1bd+v_1^2ac+v_1wbc\\
(-\delta_\alpha vv_1)C&=\ldots\quad\text{(a bore)},
&(-\delta_\alpha vv_1)D&=-vv_1d^2+v_1wcd+v_1^2c^2
\end{aligned}\right. \tag{54}
\]
\margin{N.B. the formulas for $B,D$ \unsure{simplify} by $v_1$}

\scanpage{307}

and we likewise find
\[
\eta=\unsure{(\beta,\sigma)|[\beta,\sigma]}=
\begin{pmatrix}R&S\\T&U\end{pmatrix}
\]
with
\[
\tag{55}
\begin{array}{ll}
(-vv_1\delta_\beta)R=\ldots
&(-vv_1\delta_\beta)S
 =(v_1v)su+v_1^{2}tr-v_1w'st\\
(-vv_1\delta_\beta)T=\ldots
&(-vv_1\delta_\beta)U
 =-(v_1v)u^2-v_1w'tu+v_1^{2}t^2.
\end{array}
\]
\margin{\uline{N.B.} we can divide by $v_1$.}

\unsure{We} deduce from this in particular
\deleted{\illegible{1}}
\[
\tag{56}
\left\{
\begin{aligned}
D&=\frac{-v\deleted{\unsure{v_1}}d^2+
\deleted{\unsure{v_1}}wcd+v_1^{\deleted{2}}c^2}
{-v\deleted{\unsure{v_1}}(ad-bc)}\\
\mu=1/D&=\frac{-v\deleted{\unsure{v_1}}(ad-bc)}
{-v\deleted{\unsure{v_1}}d^2+
\deleted{\unsure{v_1}}wcd+v_1^{\deleted{2}}c^2}\\
\deleted{(57)}\quad\lambda&=\deleted{\unsure{C/D}}\inserted{B/D}
=\frac{-v\deleted{\unsure{v_1}}bd+v_1^{\deleted{2}}ac+
\deleted{\unsure{v_1}}wbc}
{-v\deleted{\unsure{v_1}}d^2+
\deleted{\unsure{v_1}}wcd+v_1^{\deleted{2}}c^2}.
\end{aligned}
\right.
\]
which thus makes it possible to \unsure{make explicit}
\[
\underline{u}=\begin{pmatrix}1/D=\mu&B/D=\lambda\\0&1\end{pmatrix}
\]
in terms of $a,b,c,d$ (and the parameters $u_1,v,w$
which determine $\rho$).

For \uline{$v_1=1$}, the formula for $D$ in (56), can be written
\[
\tag{57}
D=\det(d\rho+c)/-v\delta_\alpha
=\det(d\rho+c)/\det\rho\det\alpha
\]
\deleted{the condition \unsure{to} \illegible{1} on $\alpha$ ``$D$ invertible''}
\[
\tag{58}
D(\alpha)\text{ invertible}\quad\Longleftrightarrow\quad
d\rho+c\text{ invertible}.
\]
but I realize that this is a \unsure{flat} form,

\scanpage{308}
\mspage{454}

\uline{V}\qquad
\uline{Study}\inserted{\uline{(of the equation)}}\uline{ in}
$GP(1)$, \uline{and the condition} $\mathrm{Tr}\,\sigma=0$.

Up to now we have studied equation (32)
\[
\tag{*}
[\alpha,\rho]=[\beta,\unsure{\sigma}]\,
\unsure{\varepsilon_1(\mu-1)}
\]
in $G\deleted{\unsure{\ell_2}}$ ($\simeq\mathrm{Gl}(2)$), or what amounts to the
same, \deleted{in} as a relation in $\mathrm{Sl}(2)$.
Here $\alpha,\beta$ are \deleted{\illegible{1}}\inserted{variable} sections of
$G\deleted{\unsure{\ell_2}}\simeq\mathrm{Gl}(2)$, but $[\alpha,\rho]$ and
$[\beta,\sigma]$ depend only on the \unsure{corresponding} sections
$\widetilde\alpha,\widetilde\beta$ of
\[
PG\ (\simeq GP(1))\simeq
G/\underbrace{\unsure{G_m}}_{\text{Center of }G}
\simeq\mathrm{Aut}_{\unsure{1|l}}(G)
\]
(likewise they depend on $\rho,\sigma$ only
\unsure{via} $\widetilde\rho,\widetilde\sigma$).
We can therefore \unsure{consider} that it is an equation
\inserted{in \deleted{$\mathrm{Sl}(2)$} $SG\simeq\mathrm{Sl}(2)$,}
relating $\widetilde\alpha,\widetilde\beta,\mu$.

We can propose to study the corresponding equation in
$\unsure{PG}\simeq GP(1)$,
\[
\tag{1}
[\widetilde\alpha,\widetilde\rho]
 =[\widetilde\beta,\widetilde\sigma]\cdot\widetilde\varepsilon_1(\mu-1)
\]
which can also be written in $\mathrm{Gl}(2)$, in the \unsure{equivalent} form
\[
\tag{2}
[\alpha,\rho]=\varepsilon[\beta,\sigma]\,
\overset{\deleted{\unsure{N_0}}}{\varepsilon_1}(\mu-1),
\]
where $\varepsilon$ is a section of $\unsure{G_m}$,
\deleted{and in}\inserted{of} \deleted{fact} course uniquely
\unsure{determined|determined} by $\alpha,\beta,\mu$,
(and even by $\widetilde\alpha,\widetilde\beta,\mu$)).
We will necessarily have

\scanpage{309}

\[
\tag{3}
\varepsilon^2=1
\]
(take the determinant of the two \unsure{sides}). We regard (2) as an equation
relating the variables $\alpha,\beta,\mu,\varepsilon$
($\alpha,\beta$ in $\mathrm{Gl}(2)$, $\mu,\varepsilon$ in
$\unsure{\mathbb E^1|\mathbb E'}$ -- and in fact necessarily in
$\unsure{\mathbb G_m|G_m}$, and $\varepsilon$ \unsure{already} in
$\mu_2={}_2\unsure{\mathbb G_m|G_m}$), or also
$\widetilde\alpha,\widetilde\beta,\mu,\varepsilon$.
Let us put, \deleted{as in}\inserted{\unsure{by analogy with}} (II),
\[
\tag{4}
\zeta=[\alpha,\rho],\qquad
\eta=\deleted{\illegible{1}}\varepsilon
[\unsure{\uuline{\beta}},\sigma]\deleted{\unsure{\varepsilon_1 t|\text{and}}},
\]
so that (2) takes the form
\[
\tag{5}
\zeta=\eta\varepsilon_1(\mu-1)
\]
as in (II). To be able to apply (II) to this situation, we will have to
assume that $\eta=\varepsilon[\beta,\sigma]$ satisfies the
same preliminary conditions as
\[
\eta_0=[\beta,\sigma]
\]
itself, namely $\det\eta=1$
($\Longleftrightarrow\ \varepsilon^2=1$) and
\deleted{\unsure{that}}
\[
\underbrace{\mathrm{Tr}\,\eta}_{\inserted{\varepsilon\mathrm{Tr}\,\eta_0
=\varepsilon\mathrm{Tr}\,\sigma}}
=\mathrm{Tr}\,\sigma
\qquad\deleted{\unsure{for}\ \varepsilon\mathrm{Tr}\,\eta_0=}
\]
i.e.
\[
\varepsilon\mathrm{Tr}\,\sigma=\mathrm{Tr}\,\sigma.
\]

\scanpage{310}
\mspage{455}

\uline{Lemma.} Let $\sigma$ be a section of $G$ (a form of
$\mathrm{Gl}(2)$, \unsure{associated|associated} with an Azumaya algebra
$\unsure{\mathcal M|M}$) -- i.e. a section of
$\unsure{\mathcal M|M}$ such that $\det\sigma$ invertible.
\deleted{Equivalent conditions}\inserted{Let $A_\sigma$ be the scheme}
of sections $\eta$ \deleted{\illegible{1}} of $G$ satisfying the conditions
\[
\det\eta=1\quad\text{(i.e. $\eta$ in $SG$)},\qquad
\mathrm{Tr}\,\eta\circ\sigma=\mathrm{Tr}\,\sigma.
\]
Equivalent conditions a) b) c) :

a) $A_\sigma$ is stable under translations by
$\mu_2\subset\unsure{\mathbb G_m|G_m}$ :
\[
\unsure{\eta}\longmapsto\varepsilon\eta\qquad(\varepsilon^2=1)
\]

b) $\mathrm{Tr}\,\sigma=0$

\deleted{c) $\sigma^2\in\Gamma\unsure{\mathbb G_m|G_m}$ i.e.
$\widetilde\sigma^2=1$.}

\unsure{d)|c)} (Under the \unsure{conditions} of the preceding number
\inserted{\uline{$\overline{\mathrm{IV}}$} (page 452)},
with the normalizations (34) (38) for $\rho,\sigma$) We have
\[
\tag{6}
\sigma=\begin{pmatrix}0&-1\\v'&0\end{pmatrix}
=\begin{pmatrix}0&-1\\-v&0\end{pmatrix}
\]
i.e. (with the notation of loc. cit.) we have $w=1$,
\[
\tag{7}
\text{i.e.}\qquad
\rho=\begin{pmatrix}0&1\\v&1\end{pmatrix}.
\]

\inserted{d) We have $\sigma^2\in\Gamma\unsure{\mathbb G_m|G_m}$,
i.e. $\widetilde\sigma^2=1$}

\inserted{Finally, these conditions imply the following d), and
are equivalent to \unsure{it|them} \unsure{if $\sigma$}
\deleted{\unsure{when $\sigma$}} is not a homothety on
any \unsure{fiber} (\unsure{thus} i.e. $\sigma$ is regular in the sense of
Steinberg).}

\deleted{This means that}\inserted{\unsure{The equivalence}
$a)\Longleftrightarrow b)$ means \unsure{also|thus}}
that $\unsure{\mathbb E^1|\mathbb E'}{}^{\mu_2}=(0)$, which is well
\unsure{known} and immediate to verify.
The equivalence of b) and d) is immediate, remains to prove
$b)\Longleftrightarrow c)$.
\deleted{\unsure{Let us assume|Assuming} this condition \unsure{satisfied}}
\inserted{Now we have the equation \unsure{of|of}
\unsure{Hamilt-Cayley|Hamilton-Cayley}}
\[
\sigma^2-(\mathrm{Tr}\,\sigma)\sigma+\det\sigma=0
\]

\scanpage{311}

which shows that $\mathrm{Tr}\,\sigma=0\Longrightarrow\sigma=\det\sigma\cdot1$,
and conversely if $\sigma^2$ is a scalar, we will have that
$(\mathrm{Tr}\,\sigma)\sigma=\sigma^2-\det\sigma\cdot1$ is one,
which is possible (if $\sigma$ is \deleted{\unsure{not}} a scalar on no
fiber) only if $\mathrm{Tr}\,\sigma=0$.

Let us assume these conditions satisfied for what follows,
\deleted{we \unsure{find} then that the scheme $\mathcal G_{\rho,\sigma}$
of \illegible{1} $(\alpha,\beta,\mu,\varepsilon)$ satisfying}
\inserted{Let}
\[
\tag{8}
\mathcal G^*_{\rho,\sigma}\subset
G\times G\times\unsure{\mathbb E^1|\mathbb E'}
\times\unsure{\mathbb E^1|\mathbb E'}
\]
be the scheme of solutions in $(\alpha,\beta,\mu,\varepsilon)$ of equation (2),
which is also the inverse image in
$G\times G\times\unsure{\mathbb E^1|\mathbb E'}
\times\unsure{\mathbb E^1|\mathbb E'}$ \unsure{of a} subscheme
\deleted{$\widetilde{\mathcal G}_{\rho,\unsure{\sigma}}$
\illegible{1} $PG\times$}
\[
\tag{9}
\widetilde{\mathcal G}^*_{\rho,\sigma}\subset
PG\times PG\times\unsure{\mathbb E^1|\mathbb E'}
\deleted{\times\unsure{\mathbb E^1|\mathbb E'}}
\]
consisting of the solutions \inserted{$(\alpha,\beta,\mu)$} of (1).
The projection onto the \unsure{second|last} factor
$\mathbb E^1$ gives a morphism
\[
\tag{10}
\mathcal G^*_{\rho,\sigma}\xrightarrow{\underline{\varepsilon}}\mu_2,
\]

\scanpage{312}
\mspage{456}
which moreover \unsure{factors} through $\widetilde{\mathcal G}^{*}_{\rho,\sigma}$ :
\[
\tag{11}
\widetilde{\varepsilon}:\widetilde{\mathcal G}^{*}_{\rho,\sigma}\longrightarrow\mu_2.
\]
Let us note that we have
\[
\deleted{\unsure{(11)}}\qquad
\mathcal G_{\rho,\sigma}\subset\mathcal G^{*}_{\rho,\sigma}
\]
more precisely
\[
\tag{12}
\mathcal G_{\rho,\sigma}=\underline{\varepsilon}^{-1}(1)
\]
which leads us to introduce also
\[
\tag{13}
\underset{\substack{\Vert\ \text{definition}\\
\underline{\widetilde{\varepsilon}}^{-1}(1)}}{\widetilde{\mathcal G}_{\rho,\sigma}}
=\widetilde{\mathcal G}^{*}_{\rho,\sigma},
\]
so that $\mathcal G_{\rho,\sigma}$ is the inverse image in
$G\times G\times\unsure{\mathbb E^1|\mathbb E'}$ of $\widetilde{\mathcal G}_{\rho,\sigma}$.

We have a morphism
\[
\tag{14}
\left\{\begin{aligned}
(\alpha,\beta,\mu,\varepsilon)&\longrightarrow
\left(\underbrace{(\zeta=\alpha\rho(\alpha)^{-1},\,
\eta=\varepsilon\beta\sigma(\beta)^{-1},\,\mu)}_{\in\Gamma\mathcal X_{\rho,\sigma}},\,
\underbrace{\varepsilon}_{\in\Gamma\mu_2}\right)\\
\text{from}\quad\mathcal G^{*}_{\rho,\sigma}&\longrightarrow\mathcal X_{\rho,\sigma}\times\mu_2
\end{aligned}\right.
\]
and the inverse image of a section $((\zeta,\eta,\mu),\varepsilon)$ of the right-hand side is the
scheme of solutions in $\alpha,\beta$ of the equations
\[
\tag{15}
{}^{(a)}\underset{\inserted{\unsure{\alpha}}}{\alpha\rho(\alpha)^{-1}}=\zeta,
\qquad {}^{(b)}\beta\sigma(\beta)^{-1}=\varepsilon\eta
\]

\scanpage{313}
\unsure{Now|We} we have, via the isomorphism
\[
B'_0\xrightarrow{\sim}\mathcal X_{\rho,\sigma}
\]
a privileged solution $\alpha=\underline u$ of the first equation (15), which is
\unsure{characterized} as the unique solution which is \deleted{\unsure{a}} section of
$B'_0$ (\unsure{uniqueness} resulting from
\[
T_\alpha\cap B'_0=\{1\}
\]
).
But we \unsure{will have}, for this $u$
\[
\underbrace{u\sigma(u)^{-1}}_{[u,\sigma]}=\eta
\qquad(\text{and not }\varepsilon\eta\ !)
\]
to find in terms of $u$ a solution of the equation
\inserted{(15 b)} $\beta\sigma(\beta)^{-1}=\varepsilon\eta$, we are led to try in the form
\[
\tag{16}
\beta=u\tau\quad\text{whence}\quad
\beta\deleted{\unsure{\rho}}\inserted{\sigma}(\beta)^{-1}
=\deleted{\unsure{\beta}}\inserted{u}\tau\sigma(\tau)^{-1}
\sigma(\deleted{\unsure{\beta}}\inserted{u})^{-1},
\]
trying to determine $\tau$ so that we have
\[
\tag{17}
\underbrace{\tau\sigma(\tau)^{-1}}_{[\tau,\sigma]}=\varepsilon,
\]
which indeed implies
\[
\beta\sigma(\beta)^{-1}
=\deleted{\unsure{\beta}}\inserted{u}\varepsilon\sigma(\deleted{\unsure{\beta}}\inserted{u}^{-1})
=\varepsilon\deleted{\unsure{\beta}}\inserted{u}\sigma(\deleted{\unsure{\beta}}\inserted{u})^{-1}
=\varepsilon\eta
\]
i.e. (15 b). \deleted{\unsure{Lemma}}

\scanpage{314}
\mspage{457}
The relation (17) can also be written
\[
\deleted{\sigma(\tau)=\varepsilon\tau}
\qquad(\text{take the inverse})
\]
\[
\tag{18}
\tau\sigma\tau^{-1}\sigma^{-1}=\varepsilon
\quad\text{or}\quad
\sigma\tau\sigma^{-1}\tau^{-1}=\varepsilon
\quad\text{or}
\]
\[
\tag{19}
\sigma=\varepsilon\tau(\sigma)
\quad\text{or}\quad
\tau(\sigma)=\varepsilon\sigma
\]
which implies that $\tau$ \uline{normalizes} $T_\sigma$
\[
\tag{20}
\tau(T_\sigma)=T_\sigma\quad\text{i.e.}\quad\tau\in\Gamma N_\sigma
\]
where
\[
\tag{21}
N_\sigma\overset{\text{definition}}{=}\operatorname{Norm}_G(T_\sigma).
\]

\uline{Corollary to the lemma} (page 455) the \inserted{\unsure{equivalent}} conditions of the lemma
\deleted{\unsure{are equivalent}} \inserted{are implied}
\deleted{to the} \inserted{by the} following ;
\inserted{and \unsure{imply it} if $\sigma$ is \unsure{regular|regular} :}

(e) $\forall$ \inserted{fppf local} section
$\varepsilon$ of $\mu_2$, exists \inserted{\unsure{local\textsuperscript{ly}|locally}}
section \deleted{\illegible{1}} $\tau$ of $G$ such that we have (17) (or again (19))

(f) (If $2$ invertible on $S$) idem, with $\varepsilon=-1$.

If we have (e), then we see that for every \deleted{\illegible{1}} $\varepsilon$,
we must have $\operatorname{Tr}\varepsilon\sigma=\operatorname{Tr}\sigma$
now $\operatorname{Tr}\varepsilon\sigma=\varepsilon\operatorname{Tr}\sigma$, hence
\inserted{(varying $\varepsilon$)}
$\operatorname{Tr}\sigma\in\Gamma(\unsure{\mathbb E^1|\mathbb E'}^{\mu_2}=0)$
hence $\operatorname{Tr}\sigma=0$. If we assume it only for $\varepsilon=-1$,
we will have $\operatorname{Tr}\sigma=-\operatorname{Tr}\sigma$ i.e.
$2\operatorname{Tr}\sigma=0$, i.e. $\operatorname{Tr}\sigma=0$ if $2$ invertible.
Conversely, let us assume $\operatorname{Tr}\sigma=0$, then
$\operatorname{Tr}\varepsilon\sigma=0$, moreover
$\det(\varepsilon\deleted{\unsure{\sigma}})=\varepsilon^2=1$, hence $\varepsilon$ satisfies
the ``numerical conditions'' to be of the form $[\tau,\sigma]$ for
suitable $\tau$. For

\scanpage{315}
\unsure{ensure} the local existence of $\tau$, it remains to make sure that
$\varepsilon\sigma$ is a regular section in the sense of Steinberg i.e. is nowhere
a scalar, which follows from the fact that this is so for $\sigma$.

Thus, under the present conditions, let $N'_\sigma\subset N_\sigma$ be the subscheme of
\deleted{$N_\sigma$} $G$ (or of $\unsure{N_\sigma|N}$) defined by
\[
\tag{22}
N'_\sigma=\{\tau\mid[\tau,\sigma]\in\Gamma\mu_2\}
\]
then $N'_\sigma$ is stable under right translation by $T_\sigma$ --- better,
it is a subgroup scheme of $N_\sigma$ \inserted{(containing $T_\sigma$)}, for if
$\tau,\tau'$ are such that
\[
\tau(\sigma)=\varepsilon\sigma,\qquad\tau'\sigma=\varepsilon'\sigma,
\]
we will have
\[
\tau'\tau(\sigma)=\varepsilon'\varepsilon\cdot\sigma,
\]
at the same time we see that
\[
\tau\longmapsto\varepsilon
\]
is a group homomorphism $N'_\sigma\longrightarrow\mu_2$, and we will have an exact sequence of
group schemes
\[
\tag{23}
1\longrightarrow T_\sigma\hookrightarrow N'_\sigma
\xrightarrow{\varepsilon_\sigma}\mu_2\longrightarrow1.
\]
We \unsure{will note} that if $\tau\in\Gamma G$ normalizes $T_\sigma$, we will have
\[
\tau(\sigma)=a\sigma+b\qquad a,b\in\Gamma\mathcal O_S,
\]

\scanpage{316}
\mspage{\unsure{458|488}}
and taking the traces of the two sides we find (taking into account that
$\operatorname{Tr}\tau(\sigma)=\operatorname{Tr}\sigma=0$)
$\operatorname{Tr}b.1=0$ i.e. $2b=0$, so if $b$ is invertible
$b=0$, hence $\tau(\sigma)=a\sigma$ and we necessarily have $a^2=1$ i.e. $a$ is an
$\varepsilon$. Thus
\[
\tag{24}
N'_\sigma=N_\sigma\quad\text{in residual characteristic }\ne2.
\]

\deleted{\unsure{But in characteristic}}
More \unsure{generally}, over a base $S$ such that
$2$ \deleted{not} be \deleted{\unsure{a zero divisor}} \inserted{a regular section}
(i.e. not a zero divisor) in $\Gamma(S,\mathcal O_S)$, we see that every section of
$\unsure{N'_\sigma|N_\sigma}$ is in fact a section of $N_\sigma$.
But this is no longer so in characteristic $2$, for
\[
\sigma'=a\sigma+b\qquad(a\in\Gamma G_m,\ b\in\Gamma\mathcal O_S)
\]
as soon as it satisfies
\[
\operatorname{Tr}\sigma'=\unsure{\sigma|0},\quad\det\sigma'=\det\sigma
\]
is conjugate \unsure{local\textsuperscript{ly}} to $\sigma$ --- and
\deleted{\unsure{it suffices for}} \inserted{the conditions}
are written (since
$\unsure{\det}(a\sigma+b)=a^2\det\sigma+
 ab\underbrace{\deleted{\unsure{\det}}\inserted{\operatorname{Tr}}\sigma}_{0}+b^2$
\[
\left\{\begin{aligned}
-va^2+b^2&=-v\\
2b&=0\qquad(\text{automatic in characteristic }2)
\end{aligned}\right.
\]
which in no way imply $b=0$ ....

\deleted{\illegible{1}} Let us now consider the product group
\[
B'_0\times A_\rho\times N'_\sigma,
\]

\scanpage{317}

and the morphism of schemes
\[
\tag{25}
\left\{\begin{aligned}
B'_0\times A_\rho\times N'_\sigma
&\longrightarrow \deleted{\mathcal G_{\rho,\sigma}}
 G\times G\times\unsure{\mathbb E^1|\mathbb E'}\times\unsure{\mathbb E^1|\mathbb E'}\\
(\underline u,\underline a,\underline b)
&\longmapsto
(\alpha=\underline u\underline a^{-1},\beta=\underline u\underline b^{-1},
\mu=\det\underline u,\varepsilon=\varepsilon_\sigma(\underline b)),
\end{aligned}\right.
\]

I say that the image lies in $\mathcal G^*_{\rho,\sigma}$, i.e.
that we will indeed have
\[
\tag{\unsure{*}}
\alpha\rho(\alpha)^{-1}
=\varepsilon\bigl(\beta\sigma(\beta)^{-1}\bigr)\varepsilon_1(\mu-1).
\]
Indeed
\[
\begin{aligned}
\alpha\rho(\alpha)^{-1}&=\underline u\rho(\underline u)^{-1}\\
\beta\rho(\beta)^{-1}
&=\underline u\underbrace{
\underline b^{\inserted{-1}}\sigma(\underline b^{-1})^{-1}
}_{\varepsilon_\sigma(\underline b)^{-1}=\varepsilon_\sigma(\underline b)}
\sigma(\underline u)^{-1}
=\varepsilon\underline u\sigma(\underline u)^{-1}
\end{aligned}
\]
and we know (by II) that
\[
\underline u\rho(\underline u)^{-1}
=\underline u\sigma(\underline u)^{-1}\varepsilon_1(\mu-1)
\]
hence $\unsure{(*)}$, since $\varepsilon^2=1$. But
moreover the morphism (25) is a monomorphism, \deleted{therefore}
for in terms of $(\alpha,\beta,\mu,\varepsilon)$, $u$ is determined
\inserted{(in terms of $\alpha$ already)},
(as the \uline{unique} section of $B'_0$ which
is of the form $\alpha\underline a^{-1}$, with $\underline a\in\Gamma T_\rho$
(for $T_\rho\cap B'_0=\{1\}$)), and $\underline b$ is then determined
by the relation $\beta=\underline u\underline b^{-1}$. We thus find
that (25) is an \uline{isomorphism}
\[
\tag{26}
\left\{\begin{aligned}
B'_0\times A_\rho\times N'_\sigma
&\overset{\sim}{\longrightarrow}\mathcal G^*_{\rho,\sigma}\\
(\underline u,\underline a,\underline b)
&\longmapsto
(\alpha=\underline u\underline a^{-1},\beta=\underline u\underline b^{-1},
\mu=\det\underline u,\varepsilon=\varepsilon_\sigma(\underline b))
\end{aligned}\right.
\]

\scanpage{318}
\mspage{459}

If we have two sections $(\underline u,\underline a,\underline b)$ and
$(\underline u',\underline a',\underline b')$ of
$B'_0\times T_\rho\times N'_\sigma$, under what
conditions do we have
\margin{\unsure{(27)}}
\[
\tag{27}
\widetilde\alpha=\widetilde\alpha',\quad
\widetilde\beta=\widetilde\beta'\ ?
\]
It clearly suffices to have
\[
\tag{28}
\underline a'=\lambda_1\underline a,\qquad
\underline b'=\lambda_2\underline b
\qquad \lambda_1,\lambda_2\in\Gamma G_{\mathrm m}
\]
and I say that these conditions are also necessary.
This comes from the fact that
$B'_0\hookrightarrow\unsure{PG}\simeq\unsure{GP(1)}$
is a monomorphism (let $\widetilde B'_0$ be the image of $B'_0$),
that \deleted{\unsure{the}} $A_\rho\to PG$ and
$N'_\sigma\longrightarrow PG$ have kernel $G_{\mathrm m}$
(\unsure{let|let} $\widetilde T_\rho$ and $\widetilde N'_\sigma$ be their images),
and that we have in $PG$
\[
\tag{29}
\widetilde T_\rho\cap\widetilde B'_0=\{1\},
\]
so that \deleted{\illegible{1}} $\widetilde{\underline u}$
\inserted{$\in\Gamma\widetilde B'_0$ and $\widetilde{\underline a}\in\Gamma\widetilde A_\rho$}
are defined \uline{uniquely} in terms of $\widetilde\alpha$
by the relation
\[
\deleted{\illegible{1}}\qquad
\widetilde\alpha=\widetilde{\underline u}\widetilde{\underline a}^{-1}
\]
-- and that then $\widetilde{\underline b}$ is defined
\inserted{in terms of $\widetilde\alpha,\widetilde\beta$}
by the relation
\[
\widetilde\beta=\widetilde{\underline u}\widetilde{\underline b}^{-1}.
\]

This shows that \deleted{\unsure{\(\mathcal G\)}} we have an \uline{isomorphism}
\[
\tag{30}
\left\{\begin{aligned}
\widetilde B'_0\times\widetilde A_\rho\times\widetilde N'_\sigma
&\overset{\sim}{\longrightarrow}
\widetilde{\mathcal G}^{*}_{\rho,\sigma}
\subset PG\times PG\times\unsure{\mathbb E^1|\mathbb E'}\\
(\widetilde{\underline u},\widetilde{\underline a},\widetilde{\underline b})
&\longmapsto
(\widetilde{\underline u}\widetilde{\underline a}^{-1},
\widetilde{\underline u}\widetilde{\underline b}^{-1},
\delta(\widetilde{\underline u}))
\end{aligned}\right.
\]

\scanpage{319}

where
\[
\tag{31}
\delta(\widetilde{\underline u})=\det(\underline u)
\]
when $\widetilde{\underline u}$ comes from $\underline u\in\Gamma B'_0$ --
or, what amounts to the same (cf. page 441)
$\delta(\widetilde{\underline u})$ is defined by
\[
\tag{32}
\widetilde{\underline u}(N_0)=\delta(\widetilde{\underline u})N_0.
\]

\margin{\uline{N.B.} The \unsure{centralizer} of $\widetilde\sigma$ in
$\widetilde G=PG$ \uline{is not} $\widetilde T_\sigma$, but
$\widetilde N'_\sigma$ \deleted{\illegible{2}}.}

\uline{Remark.} Let us note that a section $(\widetilde\alpha,\widetilde\beta,\mu)$
\deleted{\unsure{$\in\Gamma$}} of $\widetilde{\mathcal G}^{*}_{\rho,\sigma}$
\unsure{remains} \deleted{\unsure{entirely}} determined,
\inserted{\uline{modulo the action of} $\unsure{\widetilde N'_\sigma}$ \uline{on} $\beta$},
\unsure{to} \unsure{be} \unsure{precise}, in terms of $\widetilde\alpha$ --
\unsure{condition|conditions} to be satisfied by $\widetilde\alpha$, so that
there exist $\widetilde\beta,\mu$ \deleted{\illegible{1}} which
\unsure{complete} it to a section of $\widetilde{\mathcal G}^{*}_{\rho,\sigma}$,
being the following: the coefficient $D$ of
$\zeta=\alpha\rho(\alpha)^{-1}=\begin{pmatrix}A&B\\C&D\end{pmatrix}$
is invertible (which makes it possible to determine $\mu$ by $\mu=1/D$), then
we determine $\widetilde\eta$ \deleted{\illegible{1}} by
$\widetilde\eta=\widetilde\zeta\widetilde\varepsilon_1(\mu-1)$, and
finally we find \inserted{the} $\widetilde\beta$ \inserted{that correspond}
by the condition
\[
\widetilde\beta\widetilde\sigma(\widetilde\beta)^{-1}=\widetilde\eta.
\]

\margin{\uline{N.B.} $\widetilde\alpha$ depends on three parameters,
while $\widetilde{\mathcal G}^{*}_{\rho,\sigma}$
\uline{\unsure{is of|admits}} relative \unsure{dimension} $4=2+1+1$
(for $\widetilde{\mathcal G}^{*}_{\rho,\sigma}
\simeq\widetilde B'_0\times\widetilde T_\rho\times\widetilde N'_\sigma$).}

\margin{\uline{N.B.} We have lifted \unsure{locally$^{\mathrm t}$}
$\widetilde\alpha$ to $\alpha\in\Gamma G$,
$\zeta$ \unsure{does not} depend on the choice of lifting.}

Likewise, a section $(\alpha,\beta,\mu,\varepsilon)$
\inserted{of $\mathcal G^*_{\rho,\sigma}$}
is determined by \deleted{\unsure{$\alpha$ alone}}
\inserted{\deleted{\illegible{2}} $\alpha$ alone} modulo
\deleted{\illegible{2}} $N'_\sigma$ -- we first find
$\underline u,\underline a$ in terms

\scanpage{320}
\mspage{460}

of $\alpha$ as above, hence $\mu=\det\underline u$, hence
and it \deleted{suffices} \deleted{$\varepsilon\beta\sigma(\beta)^{-1}=\unsure{\alpha|\eta}$}
remains to determine $\beta$ as
$\beta=\underline u\underline b^{-1}$, with $\underline b\in\Gamma N'_\sigma$
arbitrary, $\varepsilon$ being defined then in terms of the choice of $\underline b$
as being $\varepsilon_\sigma(\underline b)$. The condition for
fixed $\unsure{\alpha_0|\alpha}$ to be completable by $(\beta,\mu,\varepsilon)$
to $(\alpha,\beta,\mu,\varepsilon)\in\Gamma\mathcal G^{\unsure{\alpha|*}}_{\rho,\sigma}$,
is always the condition that the coefficient $D$ \deleted{of $\zeta=\alpha\unsure{+}$}
\[
\tag{33}
\zeta=\alpha\rho(\alpha)^{-1}=\begin{pmatrix}A&B\\C&D\end{pmatrix}
\]
be invertible, which \inserted{makes it possible to} \unsure{will determine|determine} $\mu$ by
\[
\tag{34}
\mu=1/D.
\]

Let us put
\[
\tag{35}
\alpha=\begin{pmatrix}a&b\\c&d\end{pmatrix}
\]
and make $\alpha\rho(\alpha)^{-1}$ explicit; with
$\rho=\begin{pmatrix}0&1\\v&1\end{pmatrix}$

\margin{$\delta=\det\alpha=ad-bc$}
\[
\begin{aligned}
\alpha\rho(\alpha)^{-1}
&=\begin{pmatrix}a&b\\c&d\end{pmatrix}
\begin{pmatrix}0&1\\v&\deleted{\unsure{w}}\inserted{1}\end{pmatrix}
\frac1\delta\begin{pmatrix}d&-b\\-c&a\end{pmatrix}
\left(-\frac1v\right)
\begin{pmatrix}\deleted{\unsure{w}}\inserted{1}&-1\\-v&0\end{pmatrix}\\
&=-\frac1{\unsure{\delta}v}
\begin{pmatrix}bv&a+b\deleted{\unsure{w}}\\dv&c+d\deleted{\unsure{w}}\end{pmatrix}
\begin{pmatrix}d\deleted{\unsure{w}}+bv&-d\\-c\deleted{\unsure{w}}-av&c\end{pmatrix},
\end{aligned}
\]
hence
\[
\deleted{\unsure{=\frac1{\delta v}}}
\]
\[
\tag{36}
\left\{\begin{aligned}
(-\delta v)A&=\unsure{v}b(d+\unsure{vb|bv})\unsure{-}(a+b)(c+va)
&(-\delta v)B&=-vbd+c(a+b)\\
(-\delta v)C&=vd(d+vb)-(c+d)(c+va)
&(-\delta v)D&=-vd^2+cd+c^2,
\end{aligned}\right.
\]
the most essential quantity thus being
\[
\tag{37}
\left\{\begin{aligned}
(-\unsure{\delta}v)D&=-vd^2+cd+c^2=\det(d\rho+c),\\
\text{\deleted{\unsure{and after that|and after calculation}} where }&
\left\{\begin{aligned}
\deleted{\unsure{\alpha}}\ \delta&=\det\alpha=ad-bc,\\
v&\ \text{the \unsure{parameter|coefficient} which defines }
\rho=\begin{pmatrix}0&1\\v&1\end{pmatrix}
\end{aligned}\right.
\end{aligned}\right.
\]
